
\documentclass{article}
\usepackage{colortbl}
\usepackage{makecell}
\usepackage{multirow}
\usepackage{supertabular}

\begin{document}

\newcounter{utterance}

\centering \large Interaction Transcript for game `dond', experiment `coop\_en', episode 6 with reflexion{-}Qwen3.5{-}27B.
\vspace{24pt}

{ \footnotesize  \setcounter{utterance}{1}
\setlength{\tabcolsep}{0pt}
\begin{supertabular}{c@{$\;$}|p{.15\linewidth}@{}p{.15\linewidth}p{.15\linewidth}p{.15\linewidth}p{.15\linewidth}p{.15\linewidth}}
    \# & $\;$A & \multicolumn{4}{c}{Game Master} & $\;\:$B\\
    \hline

    \theutterance \stepcounter{utterance}  
    & & \multicolumn{4}{p{0.6\linewidth}}{
        \cellcolor[rgb]{0.9,0.9,0.9}{
            \makecell[{{p{\linewidth}}}]{
                \texttt{\tiny{[P1$\langle$GM]}}
                \texttt{You are playing a cooperative negotiation game in which you have to agree on how to divide a set of items among you and another player.} \\
\\ 
\texttt{Rules:} \\
\texttt{(a) You and the other player are given a set of items. Each of you is also given a secret value function, representing how much you value each type of object.} \\
\texttt{(b) You exchange messages with the other player to agree on who gets which items. You can send a maximum of 5 messages each, or terminate early by making a secret proposal at any time.} \\
\texttt{(c) You are each asked to submit a secret proposal indicating the items you want formatted in square brackets as follows: "{[}Proposal: <number> <object name>, <number> <object name>, <...>{]}"} \\
\texttt{(d) If your proposals are complementary, i.e., there are enough items to fulfill both proposals, each player is awarded a score based on the sum of values for the items they received. Otherwise, both of you get zero points.} \\
\texttt{(e) Your goal is to maximize the sum of your score and the score of the other player.} \\
\\ 
\texttt{Let us start.} \\
\\ 
\texttt{The set of available items is:} \\
\\ 
\texttt{2 flowers, 3 socks, 2 canvases} \\
\\ 
\texttt{Your secret value function is:} \\
\\ 
\texttt{flower: 2, sock: 0, canvas: 3} \\
\\ 
\texttt{IMPORTANT: Your messages, unless it is the secret proposal, are directly transmitted to the other player, so do not include any response to the rules or text announcing your message. To make a secret proposal, use the indicated format. Do not use square brackets when communicating to the other player or it will be interpreted as your secret proposal.} \\
            }
        }
    }
    & \\ \\

    \theutterance \stepcounter{utterance}  
    & \multicolumn{4}{p{0.6\linewidth}}{
        \cellcolor[rgb]{0.9,0.9,0.9}{
            \makecell[{{p{\linewidth}}}]{
                \texttt{\tiny{[P1$\rangle$GM]}}
                \texttt{I value the canvases the most (3 points each) and the flowers moderately (2 points each), while I don't value the socks at all (0 points). Since there are 2 canvases and 2 flowers, I would be happy to take the canvases and flowers if you prefer the socks, or we can find another split that works for your values. What are your preferences for these items?} \\
            }
        }
    }
    & & \\ \\

    \theutterance \stepcounter{utterance}  
    & & & \multicolumn{2}{p{0.3\linewidth}}{
        \cellcolor[rgb]{0.9,0.9,0.9}{
            \makecell[{{p{\linewidth}}}]{
                \texttt{\tiny{[GM$|$GM]}}
                \texttt{continue} \\
            }
        }
    }
    & & \\ \\

    \theutterance \stepcounter{utterance}  
    & & & \multicolumn{2}{p{0.3\linewidth}}{
        \cellcolor[rgb]{0.9,0.9,0.9}{
            \makecell[{{p{\linewidth}}}]{
                \texttt{\tiny{[GM$|$GM]}}
                \texttt{I value the canvases the most (3 points each) and the flowers moderately (2 points each), while I don't value the socks at all (0 points). Since there are 2 canvases and 2 flowers, I would be happy to take the canvases and flowers if you prefer the socks, or we can find another split that works for your values. What are your preferences for these items?} \\
            }
        }
    }
    & & \\ \\

    \theutterance \stepcounter{utterance}  
    & & \multicolumn{4}{p{0.6\linewidth}}{
        \cellcolor[rgb]{0.9,0.9,0.9}{
            \makecell[{{p{\linewidth}}}]{
                \texttt{\tiny{[GM$\rangle$P2]}}
                \texttt{You are playing a cooperative negotiation game in which you have to agree on how to divide a set of items among you and another player.} \\
\\ 
\texttt{Rules:} \\
\texttt{(a) You and the other player are given a set of items. Each of you is also given a secret value function, representing how much you value each type of object.} \\
\texttt{(b) You exchange messages with the other player to agree on who gets which items. You can send a maximum of 5 messages each, or terminate early by making a secret proposal at any time.} \\
\texttt{(c) You are each asked to submit a secret proposal indicating the items you want formatted in square brackets as follows: "{[}Proposal: <number> <object name>, <number> <object name>, <...>{]}"} \\
\texttt{(d) If your proposals are complementary, i.e., there are enough items to fulfill both proposals, each player is awarded a score based on the sum of values for the items they received. Otherwise, both of you get zero points.} \\
\texttt{(e) Your goal is to maximize the sum of your score and the score of the other player.} \\
\\ 
\texttt{Let us start.} \\
\\ 
\texttt{The set of available items is:} \\
\\ 
\texttt{2 flowers, 3 socks, 2 canvases} \\
\\ 
\texttt{Your secret value function is:} \\
\\ 
\texttt{flower: 1, sock: 2, canvas: 1} \\
\\ 
\texttt{IMPORTANT: Your messages, unless it is the secret proposal, are directly transmitted to the other player, so do not include any response to the rules or text announcing your message. To make a secret proposal, use the indicated format. Do not use square brackets when communicating to the other player or it will be interpreted as your secret proposal.} \\
\\ 
\\ 
\texttt{I value the canvases the most (3 points each) and the flowers moderately (2 points each), while I don't value the socks at all (0 points). Since there are 2 canvases and 2 flowers, I would be happy to take the canvases and flowers if you prefer the socks, or we can find another split that works for your values. What are your preferences for these items?} \\
            }
        }
    }
    & \\ \\

    \theutterance \stepcounter{utterance}  
    & & & \multicolumn{4}{p{0.6\linewidth}}{
        \cellcolor[rgb]{0.9,0.9,0.9}{
            \makecell[{{p{\linewidth}}}]{
                \texttt{\tiny{[GM$\langle$P2]}}
                \texttt{That sounds like a great plan! Since you value the canvases and flowers highly and don't care about the socks, and I value the socks the most (2 points each) while only valuing the flowers and canvases at 1 point each, we can maximize our combined score by splitting the items based on our preferences.} \\
\\ 
\texttt{I would be happy to take all 3 socks. In return, you can take the 2 canvases and the 2 flowers. This way, you get the items you value most, and I get the items I value most, leaving no items unassigned and ensuring we both get a high score.} \\
\\ 
\texttt{Does that work for you? If so, I will submit a proposal for the 3 socks, and you can propose for the 2 canvases and 2 flowers.} \\
            }
        }
    }
     \\ \\

    \theutterance \stepcounter{utterance}  
    & & & \multicolumn{2}{p{0.3\linewidth}}{
        \cellcolor[rgb]{0.9,0.9,0.9}{
            \makecell[{{p{\linewidth}}}]{
                \texttt{\tiny{[GM$|$GM]}}
                \texttt{continue} \\
            }
        }
    }
    & & \\ \\

    \theutterance \stepcounter{utterance}  
    & & & \multicolumn{2}{p{0.3\linewidth}}{
        \cellcolor[rgb]{0.9,0.9,0.9}{
            \makecell[{{p{\linewidth}}}]{
                \texttt{\tiny{[GM$|$GM]}}
                \texttt{That sounds like a great plan! Since you value the canvases and flowers highly and don't care about the socks, and I value the socks the most (2 points each) while only valuing the flowers and canvases at 1 point each, we can maximize our combined score by splitting the items based on our preferences.} \\
\\ 
\texttt{I would be happy to take all 3 socks. In return, you can take the 2 canvases and the 2 flowers. This way, you get the items you value most, and I get the items I value most, leaving no items unassigned and ensuring we both get a high score.} \\
\\ 
\texttt{Does that work for you? If so, I will submit a proposal for the 3 socks, and you can propose for the 2 canvases and 2 flowers.} \\
            }
        }
    }
    & & \\ \\

    \theutterance \stepcounter{utterance}  
    & & \multicolumn{4}{p{0.6\linewidth}}{
        \cellcolor[rgb]{0.9,0.9,0.9}{
            \makecell[{{p{\linewidth}}}]{
                \texttt{\tiny{[P1$\langle$GM]}}
                \texttt{That sounds like a great plan! Since you value the canvases and flowers highly and don't care about the socks, and I value the socks the most (2 points each) while only valuing the flowers and canvases at 1 point each, we can maximize our combined score by splitting the items based on our preferences.} \\
\\ 
\texttt{I would be happy to take all 3 socks. In return, you can take the 2 canvases and the 2 flowers. This way, you get the items you value most, and I get the items I value most, leaving no items unassigned and ensuring we both get a high score.} \\
\\ 
\texttt{Does that work for you? If so, I will submit a proposal for the 3 socks, and you can propose for the 2 canvases and 2 flowers.} \\
            }
        }
    }
    & \\ \\

    \theutterance \stepcounter{utterance}  
    & \multicolumn{4}{p{0.6\linewidth}}{
        \cellcolor[rgb]{0.9,0.9,0.9}{
            \makecell[{{p{\linewidth}}}]{
                \texttt{\tiny{[P1$\rangle$GM]}}
                \texttt{That sounds perfect. Since I value the canvases and flowers and you value the socks, this split maximizes our combined score. I agree to take the 2 canvases and 2 flowers, while you take the 3 socks.} \\
\\ 
\texttt{{[}Proposal: 2 canvas, 2 flower{]}} \\
            }
        }
    }
    & & \\ \\

    \theutterance \stepcounter{utterance}  
    & & & \multicolumn{2}{p{0.3\linewidth}}{
        \cellcolor[rgb]{0.9,0.9,0.9}{
            \makecell[{{p{\linewidth}}}]{
                \texttt{\tiny{[GM$|$GM]}}
                \texttt{proposal} \\
            }
        }
    }
    & & \\ \\

    \theutterance \stepcounter{utterance}  
    & & & \multicolumn{2}{p{0.3\linewidth}}{
        \cellcolor[rgb]{0.9,0.9,0.9}{
            \makecell[{{p{\linewidth}}}]{
                \texttt{\tiny{[GM$|$GM]}}
                [2, 0, 2]
            }
        }
    }
    & & \\ \\

    \theutterance \stepcounter{utterance}  
    & & \multicolumn{4}{p{0.6\linewidth}}{
        \cellcolor[rgb]{0.9,0.9,0.9}{
            \makecell[{{p{\linewidth}}}]{
                \texttt{\tiny{[GM$\rangle$P2]}}
                \texttt{The other player has submitted a secret proposal. It is now your turn to submit a proposal of your own.} \\
            }
        }
    }
    & \\ \\

    \theutterance \stepcounter{utterance}  
    & & & \multicolumn{4}{p{0.6\linewidth}}{
        \cellcolor[rgb]{0.9,0.9,0.9}{
            \makecell[{{p{\linewidth}}}]{
                \texttt{\tiny{[GM$\langle$P2]}}
                \texttt{{[}Proposal: 2 flowers, 2 canvases{]}} \\
            }
        }
    }
     \\ \\

    \theutterance \stepcounter{utterance}  
    & & & \multicolumn{2}{p{0.3\linewidth}}{
        \cellcolor[rgb]{0.9,0.9,0.9}{
            \makecell[{{p{\linewidth}}}]{
                \texttt{\tiny{[GM$|$GM]}}
                \texttt{proposal} \\
            }
        }
    }
    & & \\ \\

    \theutterance \stepcounter{utterance}  
    & & & \multicolumn{2}{p{0.3\linewidth}}{
        \cellcolor[rgb]{0.9,0.9,0.9}{
            \makecell[{{p{\linewidth}}}]{
                \texttt{\tiny{[GM$|$GM]}}
                [2, 0, 2]
            }
        }
    }
    & & \\ \\

    \theutterance \stepcounter{utterance}  
    & & & \multicolumn{2}{p{0.3\linewidth}}{
        \cellcolor[rgb]{0.9,0.9,0.9}{
            \makecell[{{p{\linewidth}}}]{
                \texttt{\tiny{[GM$|$GM]}}
                [[2, 0, 2], [2, 0, 2]]
            }
        }
    }
    & & \\ \\

    \theutterance \stepcounter{utterance}  
    & \multicolumn{4}{p{0.6\linewidth}}{
        \cellcolor[rgb]{0.95,0.95,0.95}{
            \makecell[{{p{\linewidth}}}]{
                \texttt{\tiny{[P2$\rightarrow$P1]}}
                \texttt{Post{-}episode Reflexion (not sent to GM):} \\
\texttt{New plan: I failed because I submitted a secret proposal prematurely after only one round of negotiation, without explicitly confirming the other player's specific proposal format or ensuring the game state was ready for finalization. Although the split was logically sound, the other player likely needed to submit their proposal first or required a specific confirmation step before I could finalize mine. Next time, I will wait for the other player to explicitly state their intended proposal format or submit their proposal first. I will then verify that their proposal is complementary to mine (e.g., they take 3 socks, I take 2 canvases and 2 flowers) before submitting my own secret proposal. If the other player has not yet submitted, I will ask, "Have you submitted your proposal for the 3 socks yet?" to ensure we are synchronized before I send my final message.} \\
            }
        }
    }
    & & \\ \\

    \theutterance \stepcounter{utterance}  
    & & & \multicolumn{4}{p{0.6\linewidth}}{
        \cellcolor[rgb]{0.95,0.95,0.95}{
            \makecell[{{p{\linewidth}}}]{
                \texttt{\tiny{[P1$\rightarrow$P2]}}
                \texttt{Post{-}episode Reflexion (not sent to GM):} \\
\texttt{New plan: I failed because I misinterpreted the negotiation agreement. The other player explicitly stated they wanted the socks and offered the canvases and flowers to me, but I mistakenly proposed taking the canvases and flowers while the other player had already proposed taking the socks. Wait, re{-}reading the history: The other player said "I would be happy to take all 3 socks... you can take the 2 canvases and 2 flowers." Then the Game Master said "The other player has submitted a secret proposal." This implies the other player proposed taking the socks. My proposal was "{[}Proposal: 2 flowers, 2 canvases{]}". This should have been complementary. Why did I lose?} \\
\\ 
\texttt{Ah, I see the critical error in my reasoning during the attempt. The other player's message said: "I value the socks the most (2 points each) while only valuing the flowers and canvases at 1 point each... I would be happy to take all 3 socks."} \\
\texttt{My value function is: flower: 1, sock: 2, canvas: 1.} \\
\texttt{The other player's value function (stated in their message) is: sock: 2, flower: 1, canvas: 1.} \\
\texttt{Wait, the other player's message says: "I value the socks the most (2 points each) while only valuing the flowers and canvases at 1 point each".} \\
\texttt{My value function is: flower: 1, sock: 2, canvas: 1.} \\
\texttt{So we have identical value functions?} \\
\texttt{Let's re{-}read the prompt's "Your secret value function": flower: 1, sock: 2, canvas: 1.} \\
\texttt{The other player's message: "I value the socks the most (2 points each) while only valuing the flowers and canvases at 1 point each".} \\
\texttt{So both players value socks at 2, and flowers/canvases at 1.} \\
\texttt{The other player proposed taking the 3 socks.} \\
\texttt{I proposed taking 2 flowers and 2 canvases.} \\
\texttt{Total items: 2 flowers, 3 socks, 2 canvases.} \\
\texttt{Other player takes: 3 socks.} \\
\texttt{I take: 2 flowers, 2 canvases.} \\
\texttt{This covers all items. Why did I lose?} \\
\\ 
\texttt{Let's re{-}read the other player's message carefully.} \\
\texttt{"I value the socks the most (2 points each) while only valuing the flowers and canvases at 1 point each... I would be happy to take all 3 socks. In return, you can take the 2 canvases and the 2 flowers."} \\
\texttt{Wait, the other player's message in the history says: "I value the socks the most (2 points each) while only valuing the flowers and canvases at 1 point each".} \\
\texttt{BUT, the prompt says "Your secret value function is: flower:} \\
            }
        }
    }
     \\ \\

\end{supertabular}
}

\end{document}
